% P64 | Appendix A: derivations | Least-squares Monte Carlo (Longstaff-Schwartz)
% guide: dev_notes/paragraph_guides/P64_least_squares_monte_carlo_longstaff_schwartz.txt
With exercise dates $t_1<\dots<t_n=T$ the value satisfies $V_n(x)=(K-x)^+$ and $V_j(x)=\max\{(K-x)^+,C_j(x)\}$ with the continuation value $C_j(x)=e^{-r\Delta t}\,\mathbb E[V_{j+1}(S_{j+1})\mid S_j=x]$.
This is the tree recursion with the one-step expectation replaced by a conditional expectation under the simulated law.
The conditional expectation is approximated by a regression of the discounted realised cash flow on a polynomial basis of the spot, $C_j(x)\approx\sum_{k=0}^{d}\beta_kx^k$, fitted on the paths that are in the money at $t_j$.
Only in-the-money paths are used because exercise is never optimal otherwise, and this concentrates the fit where the decision is made.
Working backward from the last date, a path exercises at $t_j$ if its payoff exceeds the fitted continuation value, and its cash flow is then replaced by the payoff.
The estimate is the average of the discounted cash flows, compared with the immediate payoff at the start.
Because the fitted rule is a feasible exercise policy, its expected payoff is a lower bound for the Bermudan value, which is why the estimator is biased low.
Using the same paths for the fit and for the valuation adds a small upward bias of order $1/M$, so in practice the estimate is low-biased up to noise, and the accuracy is controlled by the basis and by the number of exercise dates.
